\clearpage
\VerdierPhysicalPageAnchor{0149}%
\begingroup
\small
\setlength{\parskip}{0pt}%
\setlength{\emergencystretch}{1.5em}%
\setlength{\abovedisplayskip}{3pt plus 1pt minus 1pt}%
\setlength{\belowdisplayskip}{3pt plus 1pt minus 1pt}%

\VerdierSemanticAnchor{ch2-2-3-5-injectivity-p149-continuation}
\noindent tel que $uf=0$. En utilisant alors a), i), on voit qu'il existe
un diagramme commutatif :

\VerdierSemanticAnchor{ch2-2-3-5-injectivity-refinement-diagram}
\begin{center}
  \begin{tikzcd}[column sep=3.7em,row sep=2.3em]
    && {} \\
    & {} \arrow[ur,"v"] && {t'\in S_{\mathcal A}(\mathcal A\cap\mathcal B)} \\
    X \arrow[ur,"f"] && Y \arrow[ul,"s"] \arrow[uu,"{t'}"'] &
  \end{tikzcd}
\end{center}

\VerdierSemanticAnchor{ch2-2-3-5-injectivity-conclusion}
\noindent tel que $vf=0$. Il en résulte alors que
$m=Q_1(t')^{-1}Q_1(vf)=0$.

\VerdierSemanticAnchor{ch2-2-3-5-surjectivity-map}
Montrons maintenant que l'homomorphisme :
\[
  \operatorname{Hom}_{\mathcal A/(\mathcal A\cap\mathcal B)}(X,Y)
  \longrightarrow
  \operatorname{Hom}_{\mathcal D/\mathcal B}(X,Y)
\]
est surjectif.

\VerdierSemanticAnchor{ch2-2-3-5-surjectivity-roof}
Un élément $p\in\operatorname{Hom}_{\mathcal D/\mathcal B}(X,Y)$ se met
sous la forme :
\[
  p=Q_2(t)^{-1}Q_2(g),\qquad t\in S_{\mathcal D}(\mathcal B).
\]

\VerdierSemanticAnchor{ch2-2-3-5-surjectivity-refinement}
En utilisant a), i), on voit qu'un tel élément s'écrit toujours sous la
forme :
\[
  p=Q_2(s)^{-1}Q_2(g'),\qquad
  s\in S_{\mathcal A}(\mathcal A\cap\mathcal B),
\]

\VerdierSemanticAnchor{ch2-2-3-5-surjectivity-conclusion}
\noindent et, par suite, qu'il provient d'un élément de
$\operatorname{Hom}_{\mathcal A/(\mathcal A\cap\mathcal B)}(X,Y)$.

\VerdierSemanticAnchor{ch2-2-3-5-c-proof-i-implies-ii}
Il reste à démontrer l'équivalence de c). Démontrons tout d'abord que i)
$\Rightarrow$ ii).

\VerdierSemanticAnchor{ch2-2-3-5-iterated-quotient}
D'après (2.3.1), la catégorie $\mathcal D/\mathcal B$ est canoniquement
isomorphe à la catégorie
$(\mathcal D/(\mathcal A\cap\mathcal B))
/(\mathcal B/(\mathcal A\cap\mathcal B))$ et le foncteur canonique :

\VerdierSemanticAnchor{ch2-2-3-5-q-functor}
\[
  Q\colon\mathcal D/(\mathcal A\cap\mathcal B)
  \longrightarrow\mathcal D/\mathcal B
\]
est le foncteur de passage au quotient.

\VerdierSemanticAnchor{ch2-2-3-5-q-free-reduction}
Par suite, d'après la proposition (2.3.3, a)), pour démontrer l'implication
i) $\Rightarrow$ ii), il suffit de montrer que pour tout objet $X$ de
$\mathcal A$ et tout objet $B$ de $\mathcal B$,
$\operatorname{Hom}_{\mathcal D/(\mathcal A\cap\mathcal B)}(B,X)=0$.

\VerdierSemanticAnchor{ch2-2-3-5-m-representation}
Un élément $m$ de
$\operatorname{Hom}_{\mathcal D/(\mathcal A\cap\mathcal B)}(B,X)$ se met
sous la forme $Q_3(f)Q_3(s)^{-1}$, où :

\VerdierSemanticAnchor{ch2-2-3-5-q3-functor}
\[
  Q_3\colon\mathcal D\longrightarrow
  \mathcal D/(\mathcal A\cap\mathcal B)
\]
est le foncteur canonique, et où $s$ est un élément de
$S_{\mathcal D}(\mathcal A\cap\mathcal B)$, le système multiplicatif de
$\mathcal D$ associé à $\mathcal A\cap\mathcal B$ (2.2.2). On a alors un
diagramme de $\mathcal D$ :

\VerdierSemanticAnchor{ch2-2-3-5-w-diagram}
\begin{center}
  \begin{tikzcd}[column sep=4em,row sep=2.4em]
    & W \arrow[dl,"s"'] \arrow[dr,"f"] & \\
    B && X
  \end{tikzcd}
  \quad ,
\end{center}
\endgroup
